\section{Introdução}
\label{sc:introducao}

A recorrência de queimadas e incêndios florestais no território brasileiro representa uma das ameaças mais severas à continuidade operativa e à segurança física de infraestruturas estratégicas, com ênfase nas linhas de transmissão do Sistema Interligado Nacional (SIN). As faixas de servidão de alta tensão cruzam biomas de elevada combustibilidade (notadamente o Cerrado e áreas de transição amazônica), onde a radiação térmica extrema e as colunas de fumaça densa ionizada reduzem a rigidez dielétrica do ar, desencadeando arcos elétricos (\textit{flashover}), curtos-circuitos fase-terra e desligamentos forçados em cascata.

No âmbito da Iniciação Científica (IC) que precede e fundamenta este trabalho, concebeu-se e validou-se a primeira versão da plataforma analítica (denominada \textbf{SIRDA 1.0}): uma arquitetura Medalhão (camadas Bronze, Silver e Gold) orquestrada pelo Dagster sob o rigor do padrão FCIS (\textit{Functional Core, Imperative/OO Shell}), unificando dados geoespaciais de 7 bases públicas (INPE BDQueimadas, ONS, FUNAI, CNES/DATASUS, DNIT e IBGE) com mais de 4,05 milhões de registros. O SIRDA 1.0 consolidou o agrupamento espaço-temporal de focos via ST-DBSCAN, modelagem topológica da malha elétrica em grafos (NetworkX) e ingestão contínua em mensageria (Apache Kafka) para um banco analítico colunar (Apache Druid), sustentando consultas analíticas subsegundo sobre dashboards no Apache Superset \cite{yang2014druid}.

Contudo, a operação prática revelou que ferramentas de monitoramento puramente passivas e focadas em visualização humana tornam-se ineficientes diante de crises agudas de múltiplos focos simultâneos. A triagem manual de incidentes induz sobrecarga cognitiva nos operadores \cite{endsley1995toward} e impõe atrasos críticos na emissão de relatórios para as equipes de campo e concessionárias.

Nesse cenário, este Trabalho de Conclusão de Curso (TCC) propõe a evolução do projeto para o \textbf{SIRDA 2.0 (\textit{Actionable GeoAI DataOps})}, migrando de uma plataforma de monitoramento passivo para um ecossistema proativo de inteligência geoespacial ativa, dotado de: (i) modelagem física anisotrópica por vento (Rothermel / ERA5); (ii) indexação hexagonal discreta (Uber H3); (iii) motor de alarmes proativos via Webhooks com controle de reincidência (\textit{cooldown}); e (iv) servidor \textit{Model Context Protocol} (FastMCP) para consultas determinísticas em linguagem natural (\textit{Text-to-SQL}) e geração automática de relatórios situacionais (SitRep padrão ICS-209).
